% ============================================================
\section*{Fase 2: Implementación Técnica}
% ============================================================

\subsection*{Cifrado en Reposo: BCrypt con Salt Aleatorio}

La primera vulnerabilidad identificada en el Sistema IMC fue el uso de BCrypt con parámetros por
defecto (fuerza~10, $\approx$100~ms por hash) y la existencia de un \textbf{bypass en texto
plano} para el usuario administrador que omitía completamente la verificación criptográfica. El
OWASP recomienda un factor de costo que produzca al menos 100--300~ms de cómputo por intento de
verificación~\cite{owasp2023}.

\textbf{Cambio en \texttt{SecurityConfig.java} (línea~34):}

\begin{lstlisting}[language=Java, caption={Incremento de la fortaleza de BCrypt de 10 a 12}]
// ANTES: fuerza por defecto (10), ~100 ms por hash
return new BCryptPasswordEncoder();

// DESPUES: fuerza 12, ~300 ms por hash (recomendacion OWASP 2023)
// Resistencia 4x mayor contra ataques de GPU
return new BCryptPasswordEncoder(12);
\end{lstlisting}

BCrypt incorpora un salt aleatorio de 128 bits en cada operación de hash. El resultado tiene
la estructura:
\[
    \texttt{\$2a\$12\$[22\,chars\,salt][31\,chars\,hash]}
\]
El salt queda embebido en el hash resultante, eliminando la necesidad de almacenamiento separado
y garantizando que dos usuarios con la misma contraseña produzcan hashes completamente
distintos~\cite{owasp2023}.

\textbf{Cambio en \texttt{AuthService.java} (líneas~52--65):} Se eliminó el bypass de
autenticación en texto plano que omitía la verificación BCrypt para el administrador por defecto:

\begin{lstlisting}[language=Java, caption={Eliminación del bypass de autenticación en texto plano}]
// ANTES: comparacion en texto plano (vulnerabilidad critica)
if (DEFAULT_ADMIN_EMAIL.equals(request.getEmail()) &&
    DEFAULT_ADMIN_PASSWORD.equals(request.getPassword())) {
    // ...autenticacion sin BCrypt...
}

// DESPUES: toda autenticacion pasa por BCrypt sin excepcion
Authentication authentication = authenticationManager.authenticate(
    new UsernamePasswordAuthenticationToken(
        request.getEmail(), request.getPassword()
    )
);
\end{lstlisting}

\begin{figure}[h]
    \centering
    \includegraphics[width=0.85\textwidth]{bcrypt_config.png}
    \caption{Configuración de \texttt{BCryptPasswordEncoder(12)} en \texttt{SecurityConfig.java},
             línea~34.}
    \label{fig:bcrypt}
\end{figure}

\begin{figure}[h]
    \centering
    \includegraphics[width=0.85\textwidth]{auth_service.png}
    \caption{Método \texttt{login()} en \texttt{AuthService.java}: autenticación unificada a través
             de \texttt{AuthenticationManager} sin bypass en texto plano.}
    \label{fig:authservice}
\end{figure}

\subsection*{Seguridad en Tránsito: TLS 1.3}

El sistema carecía de configuración TLS; todas las comunicaciones entre el frontend React y el
backend Spring Boot se realizaban sobre HTTP plano (puerto~8080), exponiendo tokens JWT,
credenciales y datos biométricos a ataques de intercepción (\textit{man-in-the-middle}). El
RFC~9325 establece que TLS~1.2 y~1.3 son las únicas versiones aceptables para implementaciones
modernas, y que TLS~1.3 debe priorizarse por su modelo de seguridad mejorado~\cite{tlsrfc9325}.

Se creó el perfil de producción \texttt{application-prod.yml} con la siguiente configuración:

\begin{lstlisting}[language={}, caption={Configuración TLS 1.3 en \texttt{application-prod.yml}}]
server:
  port: 8443
  ssl:
    enabled: true
    enabled-protocols: TLSv1.3,TLSv1.2
    ciphers:
      - TLS_AES_256_GCM_SHA384        # TLS 1.3 con AES-256-GCM
      - TLS_CHACHA20_POLY1305_SHA256   # TLS 1.3, resistente a BEAST/POODLE
      - TLS_ECDHE_RSA_WITH_AES_256_GCM_SHA384  # TLS 1.2 con PFS
    key-store: ${SSL_KEYSTORE_PATH}
    key-store-password: ${SSL_KEYSTORE_PASSWORD}
    key-store-type: PKCS12
\end{lstlisting}

Los cipher suites seleccionados garantizan \textit{Perfect Forward Secrecy} (PFS), lo que
significa que el compromiso futuro de la clave privada del servidor no expone las sesiones
pasadas. Esto es una propiedad fundamental de TLS~1.3 según el RFC~9325~\cite{tlsrfc9325}.

\begin{figure}[h]
    \centering
    \includegraphics[width=0.85\textwidth]{application_prod.png}
    \caption{Perfil de producción \texttt{application-prod.yml} con TLS~1.3, cipher suites y
             configuración de keystore vía variables de entorno.}
    \label{fig:tls}
\end{figure}

\subsection*{Gestión de Llaves: Prohibición del Hardcoding}

La ISO/IEC~11770-3:2021 establece que las llaves criptográficas deben generarse, almacenarse y
destruirse mediante procesos controlados, fuera del código fuente~\cite{isoiec117703}. El sistema
IMC contenía dos violaciones críticas: el secreto JWT hardcodeado en \texttt{application.yml} y
las credenciales del administrador embebidas en \texttt{AuthService.java}.

\textbf{Cambio en \texttt{application.yml} (línea~32):}

\begin{lstlisting}[language={}, caption={Migración del secreto JWT a variable de entorno}]
# ANTES: secreto hardcodeado (visible en repositorio)
jwt:
  secret: mySecretKeyForJwtTokenGenerationAndValidationPurposeOnly

# DESPUES: referencia a variable de entorno
jwt:
  secret: ${JWT_SECRET}
  expiration: ${JWT_EXPIRATION_MS:86400000}
\end{lstlisting}

Adicionalmente, se implementó validación de entropía mínima al arrancar la aplicación en
\texttt{JwtTokenProvider.java}. HMAC-SHA512 requiere una clave de al menos 512~bits para operar
sin degradación de seguridad~\cite{nistir8413}:

\begin{lstlisting}[language=Java, caption={Validación de entropía del secreto JWT mediante \texttt{@PostConstruct}}]
private static final int MIN_SECRET_BYTES = 64; // 512 bits

@PostConstruct
public void validateSecretEntropy() {
    if (jwtSecret == null || jwtSecret.getBytes().length < MIN_SECRET_BYTES) {
        throw new IllegalStateException(
            "JWT_SECRET requiere minimo 64 bytes. " +
            "Generar con: openssl rand -hex 64"
        );
    }
}
\end{lstlisting}

\begin{figure}[h]
    \centering
    \begin{minipage}{0.48\textwidth}
        \centering
        \includegraphics[width=\textwidth]{jwt_provider.png}
        \caption{Validación de entropía mínima del secreto JWT en \texttt{JwtTokenProvider.java}.}
        \label{fig:jwt}
    \end{minipage}
    \hfill
    \begin{minipage}{0.48\textwidth}
        \centering
        \includegraphics[width=\textwidth]{env_example.png}
        \caption{Archivo \texttt{.env.example}: plantilla de variables de entorno sin secretos
                 funcionales en el repositorio.}
        \label{fig:env}
    \end{minipage}
\end{figure}

\noindent La jerarquía de almacenamiento de llaves por entorno es la siguiente:

\begin{itemize}
    \item \textbf{Desarrollo local:} variable de entorno del SO o archivo \texttt{.env} excluido
          del control de versiones mediante \texttt{.gitignore}.
    \item \textbf{CI/CD (GitHub Actions):} GitHub Secrets, cifrados en reposo y auditados.
    \item \textbf{Producción:} AWS Secrets Manager o HashiCorp Vault con rotación automática y
          acceso controlado por políticas IAM, siguiendo la norma
          ISO/IEC~11770-3:2021~\cite{isoiec117703}.
\end{itemize}

\textbf{Cambio en \texttt{authService.ts} (frontend, línea~35):} Se eliminó la persistencia de
la contraseña en \texttt{localStorage}:

\begin{lstlisting}[language={}, caption={Exclusión del campo \texttt{password} del objeto persistido en \texttt{localStorage} (TypeScript)}]
// ANTES: contrasena guardada en texto plano en localStorage
localStorage.setItem(CURRENT_USER_STORAGE_KEY, JSON.stringify(user));

// DESPUES: se omite el campo password antes de persistir
const { password: _pw, ...safeUser } = user as any;
localStorage.setItem(CURRENT_USER_STORAGE_KEY, JSON.stringify(safeUser));
\end{lstlisting}

\begin{figure}[h]
    \centering
    \includegraphics[width=0.85\textwidth]{auth_service_ts.png}
    \caption{Destructuring en \texttt{authService.ts}: el campo \texttt{password} se excluye del
             objeto que se persiste en \texttt{localStorage}.}
    \label{fig:frontend}
\end{figure}
